\section{Independent verification of the classical boundary and its
radial
complex}\label{independent-verification-of-the-classical-boundary-and-its-radial-complex}

24 September 2026. Proof locators GR0--GR7. This checks CB0--CB6 and the
stated scope of CB8 in \nolinkurl{CLASSICAL_BOUNDARY_LIFTING.md}. The previously
proved CW map used in CB7 was not independently reread during this
bounded verification. The additional exact sequence GR6 records the
complete effect of the factor of the radial coordinate at the boundary.

\subsection{GR0. Sources, the source object and the receiving
calculation}\label{gr0.-sources-the-source-object-and-the-receiving-calculation}

The human mathematical source is Alain Connes and Caterina Consani,
\href{https://arxiv.org/abs/1805.10501}{The Riemann--Roch strategy:
Complex lift of the Scaling Site, arXiv:1805.10501}, original author
source thecurve\_K.tex. The source passages checked directly are §5.4,
lines 1634--1719, especially Jdefn1, classorb1, projgnlim1, proetcov and
additivestruct; and §7.1, lines 2532--2624, especially holom, holombis,
holom1, holom2 and functionq. The original operator and moving
coefficients, including their radial factors, are retained below.

The complete user passages USR-9ad1c0a2d09dba92, USR-f55d16feb948d8c2,
USR-6152e3bc6302258c and USR-4322be19bff532cd were consulted again
together with \nolinkurl{CORPUS_AND_OPERATION_RULES.md} and the previously read
U01--U18. The receiving arithmetic below is the arithmetic already
reconstructed in the programme. None of the coordinate calculations
constructs that arithmetic from a finite sample. Primitive \(Z_1/\tau\)
receives no addition, parity, count, coordinate or numerical weight. Its
notation is not replaced by the source paper's Teichmüller notation.

Every construction below has its domain and topology specified. The
failure of one receiving function to extend to one collapsed receiving
point is not an assertion that the user's complete global construction
fails. GR2--GR3 construct the exact graph and algebra retaining its lost
directions.

\subsection{GR1. Complete classical coordinates and gauge
comparison}\label{gr1.-complete-classical-coordinates-and-gauge-comparison}

Write \(\mathbb A=\mathbb A_f\times\mathbb R\) and
\(G=\mathbb A/\mathbb Q\), with the diagonal copy of \(\mathbb Q\). The
finite coordinates admitted here are finite ideles: \[
\mathbb A_f^\times=\mathbb Q_{>0}\widehat{\mathbb Z}^{\times}.
\] This means that each coordinate is nonzero and that its valuation is
zero outside a finite set of primes. It is stronger than merely being a
finite adele with all components nonzero.

For every \(Y_f\in\mathbb A_f^\times\) there is a unique decomposition
\[
Y_f=c u,\qquad
c=\prod_p p^{v_p(Y_p)}\in\mathbb Q_{>0},\qquad
u\in\widehat{\mathbb Z}^{\times}.
\tag{GR1.1}
\] All factors occur in the displayed finite product. Their valuations
show existence. If \(c u=c'u'\), then the positive rational \(c/c'\) has
zero valuation at every prime. Its numerator and denominator therefore
have no prime factors, so \(c/c'=1\), which gives uniqueness.

On the cover with \(Y_\infty>0\), set \[
\kappa(X_f,Y_f,X_\infty,Y_\infty)
=\left([(X_f/c,X_\infty/c)],\,Y_\infty/c\right)
 = (g,y).
\tag{GR1.2}
\] For the source left action \((X,Y)\mapsto(aX+b,aY)\), where
\(a\in\mathbb Q_{>0}\) and \(b\in\mathbb Q\), the new factor is
\(c'=ac\), and hence \[
g'=[X/c+b/(ac)]=g,\qquad y'=y.
\tag{GR1.3}
\] The equality in \(G\) follows from the actual rational diagonal
translation \(b/(ac)\). A right finite unit changes \(u\) and leaves
\(c,g,y\) unchanged.

Conversely, two such points with equal \(g,y\) become equivalent after
this finite unit quotient. Indeed, take \(a=c'/c\). Equal \(y\) gives
\(Y'_\infty=aY_\infty\), and equality of \(g\) gives \[
X'/c'-X/c=b_0\in\mathbb Q.
\] Thus \(X'=aX+c'b_0\), which is the required rational affine action;
the finite second-coordinate units are matched by the separate unit
quotient. This proves the quotient-set comparison, without inferring a
topology from it.

If both nonzero signs of \(Y_\infty\) are included and the left group is
\(\mathbb Q^\times\), one may select the positive archimedean
representative explicitly by \[
\sigma=\operatorname{sgn}(Y_\infty),\qquad
g=[\sigma X/c],\qquad y=|Y_\infty|/c.
\] For arbitrary nonzero \(a\), \(c'=|a|c\) and
\(\sigma'=\sigma\operatorname{sgn}(a)\). Therefore \[
g'=g+\left[\frac{\sigma' b}{|a|c}\right]=g,\qquad y'=y.
\tag{GR1.4}
\] This sign formula is not an identification of a finite-adele
collapsed point with the archimedean zero locus of the source's separate
extension.

There is an essential topology issue. Let \(c_n=1+n!\). Then \(c_n\to1\)
in \(\mathbb A_f\): every \(c_n\) is integral at every prime, and \(n!\)
has arbitrarily large valuation at every fixed finite collection of
primes. The points \[
(X_f,Y_f,X_\infty,Y_\infty)=(0,c_n,0,1)
\] converge in the ambient additive-adele topology to \((0,1,0,1)\), but
their \(y\)-coordinates tend to zero, whereas the limit point has
\(y=1\). Consequently \(\kappa\) is not continuous for that inherited
topology, even on this classical locus. The product topology of the
source's \(G\times\mathbb R_{>0}\) chart must be carried separately. The
closure below is an explicit graph correspondence; no blowup or original
coarse compactification is inferred.

\subsection{GR2. The graph fibre over the collapsed finite-adele
boundary}\label{gr2.-the-graph-fibre-over-the-collapsed-finite-adele-boundary}

Let \[
E=\mathbb A_f^2\times\mathbb R_X\times\mathbb R_{>0,Y},
\] with the indicated product topology, and let \(\mathscr G\) be the
closure in \(E\times G\times[0,\infty)\) of the graph of (GR1.2) on the
finite-idele locus. Fix \[
e_*=(0,0,s,t),\qquad s\in\mathbb R,\quad t>0.
\] Then \[
\mathscr G_{e_*}=G\times\{0\}.
\tag{GR2.1}
\]

For the inclusion from left to right, take a convergent graph net over
\(e_*\). Eventually \(Y_f\) belongs to the open compact subgroup
\(\widehat{\mathbb Z}\). Formula (GR1.1) then makes \(c\) a positive
integer. Since \(Y_2\to0\), for every \(M\) eventually \(v_2(c)\ge M\).
It follows that \(c\ge2^M\), so \(c\to+\infty\) in the real topology.
Because \(Y_\infty\to t\), the quotient \(y=Y_\infty/c\) tends to zero.

For the reverse inclusion, first fix \(t_f\in\mathbb A_f\). There is an
integer \(d>0\) with \(dt_f\in\widehat{\mathbb Z}\), since only finitely
many negative valuations need to be cleared. Put \(c_n=d n!\). Then
\(c_n\to0\) and \(c_nt_f\to0\) in \(\mathbb A_f\), while
\(c_n\to+\infty\) in \(\mathbb R\). The graph points arising from \[
(X_f,Y_f,X_\infty,Y_\infty)=(c_nt_f,c_n,s,t)
\] approach \[
(e_*,[(t_f,0)],0).
\] Finally the image of \(\mathbb A_f\times\{0\}\) is dense in \(G\). To
verify this directly, take any \([(t_f,t_\infty)]\) and rationals
\(b_n\to t_\infty\) in \(\mathbb R\). In the quotient, \[
[(t_f-b_n,0)]-[(t_f,t_\infty)]
=[(0,b_n-t_\infty)]\longrightarrow0.
\] The fibre of the closed set \(\mathscr G\) is closed, so it contains
the closure \(G\times\{0\}\) of the points already obtained. This proves
(GR2.1).

The projection of the graph closure to \(G\times[0,\infty)\) is
continuous by construction. Its subsequent map into the source space
\((G\times\mathbb R)/\{\pm1\}\) is also continuous. Over \(e_*\) it
retains the full map \(G\to G/\{\pm1\}\), rather than identifying the
direction fibre with a single point.

\subsection{GR3. Exact restriction of the full finite coefficient
algebra}\label{gr3.-exact-restriction-of-the-full-finite-coefficient-algebra}

Retain the source algebraic coefficient ring \[
W_{\mathrm{alg}}=\mathbb C[\mathbb R_{>0}^{\times}],\qquad
\chi_\lambda([x])=x^\lambda,\qquad
\theta_\mu([x])=[x^\mu],\qquad \lambda,\mu>0.
\] For \(r\in\mathbb Q\), the source function is \[
q^r(g,y)=e_r(g)[e^{-2\pi r y}],\qquad
\chi_\lambda(q^r)=e_r(g)e^{-2\pi r\lambda y}.
\tag{GR3.1}
\] Negative and zero frequencies are included. The base of every
Teichmüller coefficient remains positive for every real \(y\).

Let \(\mathcal A\) consist of finite sums \(\sum_r w_rq^r\). Every
character evaluation extends on \(\mathscr G\) by the right side of \[
\chi_\lambda\!\left(\sum_r w_rq^r\right)
=\sum_r\chi_\lambda(w_r)e_r(g)e^{-2\pi r\lambda y}.
\tag{GR3.2}
\] This is a statement about the specified family of scalar characters,
and does not presume an unspecified topology on the group algebra.

Let \(\nu\) be Haar measure on \(G\) of any retained total mass \(M>0\).
Translation invariance gives \(\int_Ge_r\,d\nu=0\) for \(r\ne0\):
translating by an element on which \(e_r\ne1\) multiplies the integral
by a factor different from one. Thus \[
\frac1M\int_G e_r(g)\overline{e_s(g)}\,d\nu(g)
=\begin{cases}1&r=s,\\0&r\ne s.\end{cases}
\tag{GR3.3}
\] This formula retains the measure factor. Taking the source's
probability Haar measure sets \(M=1\).

The family of coefficient characters is faithful. For distinct
\(x_1,\ldots,x_d>0\), evaluation of \(\sum_ja_j[x_j]\) at
\(\lambda=1,\ldots,d\) has determinant \[
\left(\prod_jx_j\right)\prod_{i<j}(x_j-x_i)\ne0.
\] Consequently (GR3.3) followed by these character evaluations proves
that the \(\{q^r\}\) are linearly independent over \(W_{\mathrm{alg}}\).
Both \(\mathcal A\) and its generated boundary algebra are exactly
copies of \(W_{\mathrm{alg}}[\mathbb Q]\), and restriction is the
algebra isomorphism \[
\mathrm{res}:\mathcal A\longrightarrow
\mathcal B:=W_{\mathrm{alg}}\otimes\mathbb C[\mathbb Q],
\qquad wq^r\longmapsto w\otimes[r].
\tag{GR3.4}
\] The faithful boundary realization of \(w\otimes[r]\) is the family
\(\chi_\lambda(w)e_r(g)\). The inverse sends \(w\otimes[r]\) to the
actual function \(wq^r\). No rational frequency has been removed.

Equations (GR2.1) and (GR3.3) also give the exact extension criterion at
the single collapsed point in \(E\). A finite sum (GR3.2) has a
characterwise continuous extension there, independent of the approach,
exactly when \(w_r=0\) for every \(r\ne0\). Indeed, all \(g\) occur in
the graph fibre. Its boundary evaluation must therefore be constant in
\(g\). Formula (GR3.3) gives \(\chi_\lambda(w_r)=0\) for every \(r\ne0\)
and every \(\lambda>0\); faithfulness gives \(w_r=0\). The converse is
constant extension. Thus the extendable subalgebra at the collapsed
point is \(W_{\mathrm{alg}}\), while the graph retains the full algebra
\(\mathcal B\).

Under arithmetic Frobenius the finite coordinates and \(g\) remain
fixed, the radial coordinate becomes \(y/\mu\), and coefficients receive
\(\theta_\mu\). In particular, \[
\theta_\mu([e^{-2\pi r y/\mu}])=[e^{-2\pi r y}],
\] so \[
\mathfrak F_\mu\!\left(\sum_r w_rq^r\right)
=\sum_r\theta_\mu(w_r)q^r.
\tag{GR3.5}
\] Restriction (GR3.4) intertwines this with \(\theta_\mu\otimes I\) on
the full boundary algebra.

\subsection{GR4. Verification of the actual radial differential and
lifting
homotopy}\label{gr4.-verification-of-the-actual-radial-differential-and-lifting-homotopy}

The source's Remark additivestruct adds the locus \(G/\{\pm1\}\) by
extending to \((G\times\mathbb R)/\{\pm1\}\). On its oriented cover use
exactly the algebra \[
\mathcal C=C^\infty(\mathbb R)\otimes W_{\mathrm{alg}}
 \otimes\mathbb C[\mathbb Q],
\] realized by \(h(y)wq^r(g,y)\). The same Fourier and character
argument as GR3 proves that this realization is injective, now pointwise
for every \(y\); linear independence of the group-algebra basis then
makes each smooth coefficient unique.

For \(L_\lambda=y(\lambda\partial_g+i\partial_y)\), the complete
calculation is \[
\begin{aligned}
L_\lambda\chi_\lambda(h wq^r)
&=\chi_\lambda(w)e_r(g)e^{-2\pi r\lambda y}
 \left(2\pi i\lambda r yh+iyh'-2\pi i r\lambda yh\right)\\
&=iyh'\chi_\lambda(w)e_r(g)e^{-2\pi r\lambda y}.
\end{aligned}
\tag{GR4.1}
\] Therefore its actual induced differential is \(d(h wq^r)=iyh'wq^r\).

Define the chain maps in both degrees by \[
R(h wq^r)=h(0)w\otimes[r],\qquad
E(w\otimes[r])=wq^r.
\] The boundary differential is zero. We have \(RE=I\), \(Rd=0\), and
\(dE=0\).

On scalar coefficients let \[
(Th)(y)=\frac1i\int_0^y\frac{h(t)-h(0)}t\,dt.
\tag{GR4.2}
\] The integrand extends smoothly at zero because it equals
\(\int_0^1h'(ut)\,du\). This equality also supplies every derivative at
zero. Consequently \(T\) is a map on the stated smooth coefficient
space, and \[
dTh=h-h(0),\qquad Tdh=h-h(0).
\tag{GR4.3}
\] The first identity follows by differentiating (GR4.2). The second
follows by substituting \(dh(t)=it h'(t)\) and integrating. These are
respectively the degree-one and degree-zero components of \[
dT+Td=I-ER.
\tag{GR4.4}
\] Thus \(R,E\) induce inverse isomorphisms in degrees zero and one,
each with \(\mathcal B\). Moreover \(T\) preserves \(\ker R\) because
\(Th(0)=0\), so the kernel complex is contracted. This proves the zero
connecting homomorphism for this actual short exact sequence of
specified complexes.

\subsection{GR5. Frobenius, the sign involution and both cochain
degrees}\label{gr5.-frobenius-the-sign-involution-and-both-cochain-degrees}

Write \[
B_\mu(h(y)wq^r)=h(y/\mu)\theta_\mu(w)q^r.
\] Direct differentiation gives \(dB_\mu=B_\mu d\). The factor \(y/\mu\)
in the differentiated coefficient occurs on both sides; no factor has
been removed. Substitution \(t=\mu u\) in (GR4.2) gives
\(TB_\mu=B_\mu T\). The maps \(E,R\) also intertwine \(B_\mu\) with
\(\theta_\mu\otimes I\).

The source's simultaneous sign involution is \[
I(h(y)wq^r)=h(-y)wq^{-r}.
\tag{GR5.1}
\] It has the same action in both degrees of the \(L_\lambda\) complex.
Indeed \(dI=Id\), because differentiation contributes a negative sign
and the reflected radial coordinate contributes the other negative sign.
Substitution \(t=-u\), including the orientation of the integral, gives
\(TI=IT\). Boundary restriction and extension commute with the
corresponding frequency reflection. Thus all the homotopy identities
restrict to invariants. Both cohomology groups on that invariant complex
are precisely \[
\mathcal B^I=
W_{\mathrm{alg}}\otimes\mathbb C[\mathbb Q]^{r\mapsto-r}.
\tag{GR5.2}
\] This is proved by restriction of the explicit maps, without assuming
exactness of an unspecified invariant functor.

\subsection{GR6. The complete boundary contribution of the radial
factor}\label{gr6.-the-complete-boundary-contribution-of-the-radial-factor}

The raw family \(D_\lambda=\lambda\partial_g+i\partial_y\) acts on the
same realized finite algebra by \[
d_{\mathrm{raw}}(h wq^r)=ih'wq^r.
\] It defines
\(\mathcal C_{\mathrm{raw}}=(\mathcal C\xrightarrow{i\partial_y}\mathcal C)\);
the original family defines
\(\mathcal C_{\mathrm{original}}=(\mathcal C\xrightarrow{iy\partial_y}\mathcal C)\).
There is the exact chain map \[
F=(F^0,F^1)=(\mathrm{id},M_y):
\mathcal C_{\mathrm{raw}}\longrightarrow
\mathcal C_{\mathrm{original}}.
\tag{GR6.1}
\] It is a chain map because \(M_y(i\partial_y)=iy\partial_y\).

Multiplication by \(y\) is injective on smooth functions: \(yh=0\)
implies \(h(y)=0\) away from zero, and continuity gives \(h(0)=0\). Its
image is exactly the smooth functions vanishing at zero, by \[
h(y)=y\int_0^1h'(uy)\,du\quad\text{when }h(0)=0.
\] These statements apply coefficientwise to the finite tensor algebra.
Its cokernel is therefore \(\mathcal B\), by evaluation at zero. There
is a short exact sequence of complexes \[
0\longrightarrow\mathcal C_{\mathrm{raw}}
\xrightarrow{(\mathrm{id},M_y)}
\mathcal C_{\mathrm{original}}
\xrightarrow{(0,R)}
\mathcal B[-1]\longrightarrow0,
\tag{GR6.2}
\] where \(\mathcal B[-1]\) denotes \(\mathcal B\) in degree one and
zero in every other degree.

The derivative \(i\partial_y\) is surjective on \(C^\infty(\mathbb R)\),
since \(g\) has smooth preimage \((1/i)\int_0^y g(t)\,dt\). Its kernel
consists exactly of constants. Hence \[
H^0(\mathcal C_{\mathrm{raw}})=\mathcal B,\qquad
H^1(\mathcal C_{\mathrm{raw}})=0,
\] whereas GR4 gives
\(H^0(\mathcal C_{\mathrm{original}})=H^1(\mathcal C_{\mathrm{original}})=\mathcal B\).
In (GR6.2), evaluation identifies the additional degree-one class with
the boundary algebra itself. Inverting the factor \(y\) at \(y=0\) would
remove exactly this contribution and would not be an isomorphism of the
stated complexes.

The cochain actions agree with this sequence. For the raw complex, \[
d_{\mathrm{raw}}B_\mu=\mu^{-1}B_\mu d_{\mathrm{raw}},
\] so its degree-zero Frobenius is \(B_\mu\) and its degree-one
Frobenius is \(\mu^{-1}B_\mu\). For the original complex both degrees
have \(B_\mu\). The intertwining equality \[
B_\mu M_y=M_y(\mu^{-1}B_\mu)
\tag{GR6.3}
\] proves equivariance of (GR6.1) in degree one. The raw sign action is
\(I\) in degree zero and \(-I\) in degree one, since
\(d_{\mathrm{raw}}I=-Id_{\mathrm{raw}}\). The original sign action is
\(I\) in both degrees, and \[
IM_y=M_y(-I).
\tag{GR6.4}
\] Thus (GR6.2) respects both Frobenius and the source sign quotient.
These degree factors cannot be transferred unchanged from the raw
operator to the original operator at the boundary.

\subsection{GR7. Review conclusion and exact
range}\label{gr7.-review-conclusion-and-exact-range}

The graph fibre statement, coarse extension criterion, injective
boundary algebra, full radial homotopy, arithmetic Frobenius
equivariance and sign equivariance in CB0--CB6 are correct for their
specified domains. The source formula \(L_\lambda=yD_\lambda\) is
essential. GR6 proves its complete extra boundary contribution and
supplies the required comparison without inverting \(y\).

This is a lifting theorem for the actual original-character differential
on the explicitly defined smooth-radial, finite-frequency coefficient
algebra. It is not a calculation of every completed \(W\)-valued
section, every adelic support stratum, or the entire tau specialization.
It neither assigns a boundary coordinate to primitive \(Z_1/\tau\) nor
supplies Deligne's numerical weight separation by itself. The full
coefficient and frequency data survive the lift. The independent review
does not recertify the earlier CW proof cited in CB7; the latter remains
a stated dependency, with its complete proof in its own source.
